\documentclass[11pt,a4paper]{article}
\usepackage[margin=2.5cm]{geometry}
\usepackage[T1]{fontenc}
\usepackage{lmodern}
\usepackage{amsmath,amssymb}
\usepackage{enumitem}
\usepackage{microtype}
\setlength{\parindent}{0pt}
\setlength{\parskip}{0.65em}
\begin{document}
\begin{center}
{\Large\bfseries MATH6310 Metric Spaces and Topology}\\[3pt]
{\large Exercise 1.7 --- Solution}
\end{center}

\textbf{(a) Define the inverse map.}

For any $f\in F$, define
\[
\boxed{\sigma(f)=\{s\in S:f(s)=1\}.}
\]
In other words, $\sigma(f)$ is the set of all elements of $S$ at which $f$ takes the value $1$.

This is a subset of $S$, so $\sigma(f)\in\mathcal P(S)=P$. Therefore $\sigma:F\to P$ is well-defined.

\textbf{(b) Show the two compositions are identity maps.}

\textbf{First identity: $\sigma\circ\chi=\operatorname{id}_P$.}

Take any $A\in P$ (so $A\subseteq S$). Applying $\chi$ first gives $\chi(A)=\mathbf 1_A$. Applying $\sigma$ next gives
\begin{align*}
(\sigma\circ\chi)(A)
  &=\sigma(\chi(A))\\
  &=\sigma(\mathbf 1_A)\\
  &=\{s\in S:\mathbf 1_A(s)=1\}\\
  &=\{s\in S:s\in A\}\\
  &=A.
\end{align*}
The fourth line follows because $\mathbf 1_A(s)=1$ exactly when $s\in A$.
Since this holds for \emph{every} $A\in P$, we have
\[
\boxed{\sigma\circ\chi=\operatorname{id}_P.}
\]

\textbf{Second identity: $\chi\circ\sigma=\operatorname{id}_F$.}

Take any $f\in F$. Put $A=\sigma(f)=\{s\in S:f(s)=1\}$.
Then
\[
(\chi\circ\sigma)(f)=\chi(A)=\mathbf 1_A.
\]
We need to check that $\mathbf 1_A=f$. Two functions with the same domain and codomain are equal if they give the same value at every input.

Let $s\in S$. There are only two cases, because $f(s)\in\{0,1\}$:
\begin{enumerate}[label=\textbf{Case \arabic*:},leftmargin=*,itemsep=1pt]
\item If $f(s)=1$, then $s\in A$ by the definition of $A$. Therefore $\mathbf 1_A(s)=1=f(s)$.
\item If $f(s)=0$, then $s\notin A$. Therefore $\mathbf 1_A(s)=0=f(s)$.
\end{enumerate}
In both cases, $\mathbf 1_A(s)=f(s)$. As this is true for every $s\in S$, we have $\mathbf 1_A=f$. Hence
\[
(\chi\circ\sigma)(f)=f
\quad\text{for every }f\in F,
\]
so
\[
\boxed{\chi\circ\sigma=\operatorname{id}_F.}
\]

\textbf{Conclusion.} We have shown that $\sigma$ reverses $\chi$ in both directions. Thus $\sigma$ is the inverse of $\chi$, and therefore $\chi$ is a bijection:
\[
\boxed{\mathcal P(S)\cong\mathcal F(S;\{0,1\}).}\qquad\square
\]

\end{document}
